\subsection{Linear Operators as Modules}

\begin{definition} \label{def:operator-invariant-factors-elementary-divisors} \hlt{Invariant Factors and Elementary Divisors}\\
Let $T\in\LL(\mathcal{V})$.
\begin{enumerate}[label=(\roman*)]
\setlength{\itemsep}{0pt}
\item A \hlt{cyclic decomposition} of $T$ is a direct sum $\mathcal{V}=\bigoplus\limits_{j=1}^r\mathcal{C}_j$ of nonzero $T$-invariant cyclic subspaces.
\item If $\mathcal{C}_j\cong\mathbb{F}[z]/\langle f_j\rangle$, $f_j\in\mathbb{F}[z]$ monic and nonconstant, $f_r\mid f_{r-1}\mid\cdots\mid f_1$,\\
then $f_1,\ldots,f_r$ are the \hlt{invariant factors} of $T$.
\item Prime-power factors in factorizations of $f_j$ into monic irreducibles are the \hlt{elementary divisors} of $T$.
\end{enumerate}
\end{definition}

\begin{theorem} \label{thm:operator-invariant-factor-skeleton} \hlt{Invariant-Factor Skeleton for Operators}\\
Let $\mathcal{V}\neq\{\mathbf{0}\}$ be finite-dimensional over $\mathbb{F}$, and let $T\in\LL(\mathcal{V})$. Equip $\mathcal{V}$ with $\mathbb{F}[z]$-module structure $p\cdot\mathbf{v}=p(T)\mathbf{v}$.\\
Then $\exists$ monic nonconstant polynomials $f_1,\ldots,f_r\in\mathbb{F}[z]$ with $f_r\mid f_{r-1}\mid\cdots\mid f_1$ and nonzero $T$-invariant cyclic subspaces $\mathcal{C}_1,\ldots,\mathcal{C}_r$ such that $\mathcal{V}=\mathcal{C}_1\oplus\cdots\oplus\mathcal{C}_r$ and $\mathcal{C}_j\cong \mathbb{F}[z]/\langle f_j\rangle$.\\
In the concatenated cyclic bases, $[T]=C_{f_1}\oplus\cdots\oplus C_{f_r}$. Moreover, $f_1=\mu_T$, the characteristic polynomial of $T$ is $f_1\cdots f_r$, and the polynomials $f_1,\ldots,f_r$ are uniquely determined by $T$.
\begin{proof}
As an $\mathbb{F}[z]$-module, $\mathcal{V}$ is finitely generated, and $\mu_T(T)=0$ makes every vector torsion.\\
Invariant-factor theorem over PID $\mathbb{F}[z]$ gives $\mathcal{V}\cong\bigoplus\limits_{j=1}^r\mathbb{F}[z]/\langle f_j\rangle$, $f_r\mid\cdots\mid f_1$, with each $f_j$ monic and nonconstant.\\
Transporting cyclic generators gives subspaces $\mathcal{C}_j$, and Theorem \ref{thm:cyclic-subspace-quotient-and-companion} gives companion block $C_{f_j}$ on each cyclic basis. The annihilator of the direct sum is $\langle\operatorname{lcm}(f_1,\ldots,f_r)\rangle=\langle f_1\rangle$, while the annihilator of the operator module is $\langle\mu_T\rangle$; hence $f_1=\mu_T$. Also $p_{C_{f_j}}=f_j$ by Theorem \ref{thm:companion-matrix-has-prescribed-minimal-polynomial}, so the block diagonal form has characteristic polynomial $\prod_{j=1}^r f_j$. Uniqueness is exactly uniqueness of invariant factors in the PID structure theorem.\\
\hlt{Structural Commentary}: Torsion module decomposition produces companion blocks.\\
\hlt{Key Insight}: Invariant factors are the polynomial-module version of choosing canonical blocks for similarity.
\end{proof}
\end{theorem}

\begin{perspective} \label{per:elementary-divisors-and-jordan-blocks} \hlt{Elementary Divisors and Jordan Blocks}\\
Over $\mathbb{C}$, every elementary divisor has the form $(z-\lambda)^k$. The cyclic summand $\mathbb{C}[z]/\langle (z-\lambda)^k\rangle$ corresponds to one Jordan block $J_k(\lambda)$. Hence elementary divisors over $\mathbb{C}$ are precisely Jordan blocks counted with multiplicity.
\end{perspective}

\begin{perspective} \label{per:operator-module-similarity-dictionary} \hlt{Operator Modules and Similarity}\\
For each $T\in\LL(\mathcal{V})$, let $\mathcal{V}_T$ be the $\mathbb{F}[z]$-module with action $p\cdot\mathbf{v}=p(T)\mathbf{v}$. For $S,T\in\LL(\mathcal{V})$, $S\sim T \iff \mathcal{V}_S\cong\mathcal{V}_T$, as $\mathbb{F}[z]$-modules, and each cyclic summand $\mathbb{F}[z]/\langle f\rangle$ corresponds to the companion block $C_f$ by Theorem \ref{thm:cyclic-subspace-quotient-and-companion}.
\end{perspective}

\begin{definition} \label{def:rational-canonical-block-forms} \hlt{Rational Canonical Block Forms}
\begin{enumerate}[label=(\roman*)]
\setlength{\itemsep}{0pt}
\item A matrix is in \hlt{invariant-factor rational canonical form} if $A=\operatorname{diag}(C_{s_1},\ldots,C_{s_r})$, $s_r\mid s_{r-1}\mid\cdots\mid s_1$, where each $s_j\in\mathbb{F}[z]$ is monic and nonconstant; the $s_j$ are the \hlt{invariant factors}.
\item A matrix is in \hlt{elementary-divisor rational canonical form} if $A=\operatorname{diag}(C_{q_1},\ldots,C_{q_m})$, where each $q_j$ is a power of a monic irreducible polynomial; the $q_j$ are the \hlt{elementary divisors}.
\end{enumerate}
\end{definition}

\begin{theorem} \label{thm:rational-canonical-form-skeleton} \hlt{Rational Canonical Form Skeleton}\\
Let $A\in\operatorname{Mat}(n,n,\mathbb{F})$. $\exists$ unique monic nonconstant $s_1,\ldots,s_r\in\mathbb{F}[z]$ with $s_r\mid\cdots\mid s_1$ such that $A$ is similar over $\mathbb{F}$ to $\operatorname{diag}(C_{s_1},\ldots,C_{s_r})$. Moreover, $\mu_A=s_1,\qquad p_A=s_1s_2\cdots s_r$.\\
Equivalently, the elementary divisors, counted with multiplicity, form a complete similarity invariant over $\mathbb{F}$.
\begin{proof}
Let $T_A:\mathbb{F}^n\to\mathbb{F}^n$ be $T_A\mathbf{x}=A\mathbf{x}$. By Theorem \ref{thm:operator-invariant-factor-skeleton}, $\mathbb{F}^n=\mathcal{C}_1\oplus\cdots\oplus\mathcal{C}_r$; $\mathcal{C}_j\cong\mathbb{F}[z]/\langle s_j\rangle$, $s_r\mid\cdots\mid s_1$.\\
In cyclic bases, $[T_A]=\operatorname{diag}(C_{s_1},\ldots,C_{s_r})$, so $A$ is similar to this form.\\
Uniqueness, $\mu_A=s_1$, $p_A=s_1\cdots s_r$, and the elementary-divisor classification follow from the same theorem.\\
\hlt{Structural Commentary}: Companion blocks classify similarity without requiring $p_A$ to split.\\
\hlt{Key Insight}: Similarity classes can be represented by polynomial data rather than by choosing roots.
\end{proof}
\end{theorem}

\begin{figure}[H]
  \centering
  \begin{tikzpicture}[
  >=Stealth,
  every node/.style={font=\small},
  class/.style={draw=blue!60!black, fill=blue!7, rounded corners=2pt,
    minimum width=2.25cm, minimum height=0.72cm, align=center, inner sep=5pt},
  data/.style={draw=dartmouthgreen!65!black, fill=dartmouthgreen!8, rounded corners=2pt,
    minimum width=2.9cm, minimum height=0.72cm, align=center, inner sep=5pt},
  form/.style={draw=gray!60, fill=gray!6, rounded corners=2pt,
    minimum width=3.45cm, minimum height=0.72cm, align=center, inner sep=5pt},
  link/.style={->, line width=0.75pt, draw=black!70},
  note/.style={font=\scriptsize, inner sep=2pt, fill=white}
]
  \node[class] (ops) at (0,0) {$[T]_{\sim}$};
  \node[class] (mods) at (4.3,0) {$[\mathcal{V}_T]_{\mathbb{F}[z]}$};
  \node[data] (invs) at (8.7,0) {$\{s_i\}:s_r\mid\cdots\mid s_1$\\$\{q_j\}$};
  \node[form] (mats) at (13.4,0) {$\operatorname{diag}(C_{s_1},\ldots,C_{s_r})$};

  \draw[link] (ops) -- node[note, above] {$p\cdot\mathbf{v}=p(T)\mathbf{v}$} (mods);
  \draw[link] (mods) -- node[note, above] {$\bigoplus_i\mathbb{F}[z]/\langle s_i\rangle$} (invs);
  \draw[link] (invs) -- node[note, above] {$s_i\mapsto C_{s_i}$} (mats);
\end{tikzpicture}

  \caption{Rational-canonical dictionary: $[T]_{\sim}$, $[\mathcal{V}_T]_{\mathbb{F}[z]}$, $\{s_i\}$ or $\{q_j\}$, and $\operatorname{diag}(C_{s_i})$ are equivalent data.}
  \label{fig:operator-module-canonical-form-dictionary}
\end{figure}
